\documentclass[11pt,a4paper,oneside,openany]{book}
\usepackage{biblioflow-ista}

\renewcommand{\titolodocumento}{Report di Progetto}
\renewcommand{\sottotitolodocumento}{Report del Progetto}

\begin{document}
\frontmatter
\copertina
\tableofcontents
\listoffigures

\mainmatter
\chapter{Descrizione del Sistema}

\section{Contesto Applicativo}

Nelle biblioteche universitarie i posti studio sono una risorsa limitata e
molto contesa, soprattutto nei periodi d'esame. In assenza di un sistema di
prenotazione l'accesso segue la regola del primo arrivato: chi abita vicino
alla sede occupa i posti nelle prime ore del mattino, mentre gli studenti
pendolari, che dipendono dagli orari dei mezzi pubblici, arrivano quando la sala
è già piena. È diffusa anche l'occupazione \enquote{informale} del posto,
lasciando libri o zaino per ore, che rende la disponibilità reale diversa da
quella apparente.

Il personale della biblioteca, d'altra parte, non dispone di dati affidabili
sull'occupazione delle sale, sui posti che restano vuoti pur essendo stati
promessi a qualcuno e sugli studenti che si prenotano senza presentarsi
(\emph{no-show}). Infine, gli studenti con esigenze di accessibilità hanno
bisogno di sapere in anticipo quali posti sono raggiungibili e attrezzati.

La digitalizzazione della prenotazione dei posti, collegata a una conferma
della presenza all'arrivo, permette di rendere l'accesso più equo e di liberare
automaticamente i posti non utilizzati.

\section{La Piattaforma BiblioFlow}

BiblioFlow nasce in questo contesto come piattaforma web per la gestione dei
posti studio e dei servizi di una biblioteca universitaria. È un'applicazione
web progressiva (PWA) sviluppata con Next.js, React e TypeScript: si usa dal
browser ma può essere installata su smartphone e continua a mostrare una
pagina di cortesia in assenza di rete. I dati sono conservati in PostgreSQL e
gestiti tramite l'ORM Prisma; l'autenticazione è affidata ad Auth.js
(NextAuth), con accesso tramite email e password o account Google.

Il sistema distingue tre ruoli:
\begin{itemize}
  \item \textbf{Studente:} prenota un posto scegliendolo sulla mappa della
        sala, conferma la presenza con il check-in, consulta il catalogo,
        chiede libri in prestito o la loro preparazione al banco e riceve
        notifiche;
  \item \textbf{Bibliotecario:} gestisce prenotazioni, posti, prestiti,
        richieste e anomalie, effettua il check-in degli studenti tramite lo
        scanner dei QR code e consulta le statistiche;
  \item \textbf{Amministratore:} ha gli stessi strumenti del bibliotecario e
        in più può attivare o disattivare gli account.
\end{itemize}

Due caratteristiche distinguono BiblioFlow da un semplice calendario di
prenotazione: l'opzione \textbf{pendolare}, che aggiunge un margine di
tolleranza sull'orario di arrivo per chi viaggia con i mezzi pubblici, e
l'attenzione all'\textbf{accessibilità}, con posti accessibili segnalati sulla
mappa, preferenze di visualizzazione (alto contrasto, testo grande, movimento
ridotto) e navigazione completa da tastiera.

\subsection{Analisi del Funzionamento}

Il funzionamento di BiblioFlow ruota attorno al ciclo di vita della
prenotazione di un posto e può essere descritto attraverso i seguenti
passaggi:
\begin{itemize}
  \item \textbf{Autenticazione:} lo studente accede con email e password
        (verificata con bcrypt) o con Google. La sessione è un token JWT
        conservato in un cookie; un middleware intercetta le richieste e
        respinge quelle prive del cookie di sessione, mentre i controlli sul
        ruolo sono svolti dalle singole route.
  \item \textbf{Prenotazione guidata:} una procedura in quattro passi raccoglie
        data, durata (due ore, mezza giornata di mattina o di pomeriggio,
        giornata intera), sala e posto. La scelta del posto avviene su una
        mappa interattiva che mostra lo stato di ogni posto e le sue
        caratteristiche (presa elettrica, accessibilità). Il server verifica gli
        orari di apertura della sala e l'assenza di sovrapposizioni prima di
        creare la prenotazione nello stato \texttt{CONFERMATA}.
  \item \textbf{Check-in e uso del posto:} all'arrivo lo studente conferma la
        presenza dalla home, dove la prenotazione del giorno è accompagnata da
        un QR code, oppure si fa scansionare il QR dal personale. Il check-in è
        ammesso in una finestra di quindici minuti attorno all'inizio della
        prenotazione e porta la prenotazione in \texttt{CHECK\_IN} e il posto in
        \texttt{OCCUPATO}; al termine la prenotazione diventa
        \texttt{COMPLETATA}.
  \item \textbf{Automazioni:} un'attività pianificata invia i promemoria di
        check-in, avvisa delle scadenze dei prestiti e rilascia i posti delle
        prenotazioni non confermate entro il margine (stato \texttt{NO\_SHOW}).
  \item \textbf{Notifiche e tracciabilità:} ogni operazione rilevante produce
        una notifica persistente per l'utente e un evento nel registro
        \texttt{LogEvento}, che alimenta la dashboard e le statistiche dello
        staff.
\end{itemize}

\chapter{Struttura del Sistema}

L'architettura di BiblioFlow segue la struttura di un'applicazione Next.js con
App Router: pagine e route handler convivono nella cartella \file{src/app},
mentre la logica condivisa, i componenti dell'interfaccia e il modello dei dati
sono organizzati in package distinti, ognuno preposto a una specifica
responsabilità. Di seguito viene presentata l'analisi dei moduli principali,
così come si presentano nella versione di partenza del progetto.

\section{Package Principali del Sistema}

\subsection{app (pagine)}

L'obiettivo è fornire le schermate dell'applicazione. Ogni cartella sotto
\file{src/app} corrisponde a un percorso dell'interfaccia; le pagine dello
studente sono componenti client che dialogano con le API tramite
\texttt{fetch}, mentre la dashboard dello staff è resa lato server e legge i
dati direttamente con Prisma.

Moduli principali:
\begin{itemize}
  \item \file{layout.tsx}: struttura comune a tutte le pagine; monta i provider
        di sessione, tema, accessibilità e PWA, l'intestazione e i componenti di
        accessibilità (salto al contenuto, annunci per i lettori di schermo).
  \item \file{page.tsx}: home dello studente; mostra la prenotazione attiva con
        il QR code per il check-in e le scorciatoie verso le funzioni
        principali.
  \item \file{login}, \file{registrazione}, \file{recupera-password},
        \file{reset-password}: accesso e gestione delle credenziali.
  \item \file{prenota}: procedura guidata di prenotazione in quattro passi
        (data, durata, sala, posto) con mappa interattiva e opzione pendolare.
  \item \file{prenotazioni} e \file{prenotazioni/[id]/estendi}: elenco delle
        prenotazioni dello studente, cancellazione ed estensione dell'orario.
  \item \file{libri}, \file{libri/[id]}, \file{prestiti}: catalogo, dettaglio
        di un libro, richiesta e rinnovo dei prestiti.
  \item \file{notifiche}, \file{profilo}: centro notifiche e dati del profilo,
        comprese le preferenze di pendolarismo e accessibilità.
  \item \file{admin/*}: area dello staff, con dashboard, gestione di
        prenotazioni, posti, prestiti, richieste di preparazione, utenti e
        anomalie, scanner dei QR code e statistiche.
\end{itemize}

\diagramma[0.62\linewidth]{analisi-modulo-app}{Analisi modulo app (pagine).}

\subsection{app/api (route handler)}

Questo package rappresenta il backend dell'applicazione. Ogni file
\file{route.ts} espone uno o più metodi HTTP e restituisce JSON; le route
validano l'input, verificano la sessione dove previsto e accedono ai dati
tramite Prisma.

Moduli principali:
\begin{itemize}
  \item \file{auth/*}: gestore di Auth.js e flussi di registrazione,
        verifica dell'email, recupero e reimpostazione della password, con
        limitazione della frequenza delle richieste.
  \item \file{prenotazioni}, \file{prenotazioni/[id]},
        \file{prenotazioni/[id]/check-in}, \file{prenotazioni/[id]/estendi}:
        elenco, creazione, lettura, cancellazione, check-in ed estensione delle
        prenotazioni.
  \item \file{posti}, \file{sale}, \file{disponibilita-giorni}: stato dei
        posti, sale con i relativi orari e giorni prenotabili.
  \item \file{libri}, \file{prestiti}, \file{richieste}: catalogo, prestiti e
        rinnovi, richieste di preparazione dei libri al banco.
  \item \file{notifiche}, \file{profilo}: lettura, marcatura ed eliminazione
        delle notifiche; lettura e modifica del profilo.
  \item \file{admin/*}: operazioni dello staff (annullamento e modifica delle
        prenotazioni, stato dei posti, prestiti, richieste, utenti, anomalie,
        statistiche, validazione dei QR code).
  \item \file{cron/automations}: punto di ingresso dell'attività pianificata,
        protetto da un segreto condiviso.
  \item \file{sse/posti}: canale Server-Sent Events per gli aggiornamenti in
        tempo reale dei posti.
\end{itemize}

A monte delle route opera \file{src/middleware.ts}, che lascia passare le
pagine e le API pubbliche (accesso, registrazione, \file{api/auth},
\file{api/cron}) e per tutte le altre verifica la presenza del cookie di
sessione: in sua assenza risponde 401 alle API e reindirizza le pagine al
login, conservando la destinazione richiesta.

\diagramma[0.9\linewidth]{analisi-modulo-api}{Analisi modulo app/api.}

\subsection{lib}

Raccoglie i servizi condivisi dal backend: autenticazione, accesso ai dati,
automazioni e utilità trasversali.

Moduli principali:
\begin{itemize}
  \item \file{auth.ts}: configurazione di Auth.js con i provider Credentials e
        Google, i callback che arricchiscono il token con ruolo e preferenze
        dell'utente e le funzioni \texttt{hasRole} e \texttt{isStaff}.
  \item \file{prisma.ts}: istanza unica del client Prisma collegata a
        PostgreSQL tramite l'adattatore \texttt{pg}.
  \item \file{automation-service.ts}: logica delle automazioni periodiche:
        \texttt{sendCheckInReminders}, \texttt{sendLoanExpiryAlerts},
        \texttt{releaseNoShowReservations}, \texttt{notifyAvailableSeat} e
        l'orchestratore \texttt{runAllAutomations}.
  \item \file{realtime-events.ts} e \file{sse-emitter.ts}: emettitore di
        eventi in memoria e funzioni di pubblicazione degli aggiornamenti dei
        posti e delle notifiche sul canale SSE.
  \item \file{qr-signature.ts}: firma HMAC e verifica dei QR code usati per il
        check-in, con controllo della scadenza.
  \item \file{rate-limit.ts}: limitatori di frequenza per indirizzo IP
        (accesso, registrazione, prenotazioni, prestiti, API).
  \item \file{password.ts}, \file{env.ts}, \file{utils.ts}: hashing e
        regole delle password, validazione delle variabili d'ambiente
        all'avvio, composizione delle classi CSS.
\end{itemize}

\diagramma[0.85\linewidth]{analisi-modulo-lib}{Analisi modulo lib.}

\subsection{components}

Contiene i componenti React riutilizzabili che compongono l'interfaccia.
Questo approccio modulare facilita la manutenzione e la coerenza visiva tra
l'area dello studente e quella dello staff.

Moduli principali:
\begin{itemize}
  \item \file{mappa-biblioteca.tsx}: mappa interattiva dei posti di una sala,
        con legenda degli stati e delle caratteristiche; sotto i 768 pixel di
        larghezza cede il posto a \file{mobile-posti-grid.tsx}, una griglia
        pensata per lo schermo dello smartphone.
  \item \file{qrcode-checkin.tsx}: mostra il QR code della prenotazione e
        invia la conferma di check-in.
  \item \file{admin/*}: menu di azioni su prenotazioni, posti, prestiti,
        utenti e anomalie, filtri, card della dashboard e grafici delle
        statistiche.
  \item \file{layout/*} e \file{providers/*}: intestazione con il contatore
        delle notifiche, barra laterale dello staff, provider di sessione, tema
        e PWA.
  \item \file{ui/*}: primitive grafiche (pulsanti, card, dialog, menu, schede,
        selettori) costruite su Radix UI.
\end{itemize}

\diagramma[0.85\linewidth]{analisi-modulo-components}{Analisi modulo components.}

\subsection{accessibility e contexts}

Questo package raccoglie il supporto all'accessibilità, che nel progetto è un
requisito e non un'aggiunta.
\begin{itemize}
  \item \file{contexts/accessibility-context.tsx}: stato globale delle
        preferenze (alto contrasto, testo grande, dimensione del carattere,
        movimento ridotto, tema scuro, navigazione da tastiera), salvate nel
        browser e applicate all'intera interfaccia.
  \item \file{components/accessibility/*}: collegamento per saltare al
        contenuto, scorciatoie da tastiera, intrappolamento del focus nei
        dialog e annunci per i lettori di schermo (\texttt{LiveAnnouncer}).
\end{itemize}

\subsection{hooks e types}

\begin{itemize}
  \item \file{hooks/use-sse.ts}: hook React che apre una connessione
        \texttt{EventSource}, gestisce la riconnessione e consegna gli eventi
        ricevuti; \texttt{usePostiRealtime} lo specializza per il canale dei
        posti.
  \item \file{types/index.ts}: tipi TypeScript condivisi tra pagine e
        componenti.
\end{itemize}

\subsection{prisma}

Questa cartella contiene il modello dei dati e gli strumenti per crearlo.
\begin{itemize}
  \item \file{schema.prisma}: definisce undici entità, tra cui utenti, sale,
        posti, prenotazioni, libri, prestiti, notifiche e registro degli
        eventi, insieme agli stati di posti, prenotazioni e prestiti. La data di una prenotazione è salvata come
        \texttt{@db.Date} e i suoi orari come \texttt{@db.Time}.
  \item \file{migrations/*}: storico delle migrazioni del database.
  \item \file{seed.ts}: popola un'installazione nuova con sale, posti, libri e
        account dimostrativi per ciascun ruolo.
\end{itemize}

\diagramma[0.95\linewidth]{modello-dati-as-is}{Analisi modulo prisma (modello dei dati).}

\subsection{public (PWA)}

Contiene le risorse statiche che rendono BiblioFlow installabile e
utilizzabile anche con una connessione instabile.
\begin{itemize}
  \item \file{manifest.json}: nome, icone, colori e scorciatoie
        dell'applicazione installata.
  \item \file{sw.js}: service worker con strategia \emph{network first} per le
        pagine e le API e \emph{cache first} per le risorse statiche; senza
        rete mostra \file{offline.html}.
  \item \file{components/pwa/*}: banner di installazione e indicatore dello
        stato offline.
\end{itemize}

\chapter{Flusso di Funzionamento del Sistema}

\section{Tipologie di Utilizzo}

BiblioFlow è stato progettato per accompagnare la giornata in biblioteca dal
punto di vista di chi studia, di chi gestisce la sala e del sistema stesso,
che interviene in autonomia a intervalli regolari. Il sistema offre quindi tre
modalità principali di utilizzo:
\begin{itemize}
  \item \textbf{Utilizzo da parte dello studente:} tramite browser o
        applicazione installata, lo studente prenota il posto, riceve i
        promemoria e conferma la presenza con il check-in.
  \item \textbf{Utilizzo da parte dello staff:} bibliotecari e amministratori
        usano l'area \file{/admin} per effettuare il check-in con lo scanner,
        annullare o modificare prenotazioni e risolvere le anomalie.
  \item \textbf{Automazioni pianificate:} un'attività pianificata invoca
        periodicamente l'endpoint \file{/api/cron/automations}, che esegue
        promemoria, avvisi e rilascio dei posti non occupati.
\end{itemize}

\section{Flusso di Utilizzo}

In questa sezione vengono dettagliati i flussi operativi che caratterizzano le
tre modalità di utilizzo. Per ognuna si descrive la sequenza di eventi che
porta dall'azione dell'utente, o dell'attività pianificata, al risultato
finale.

\subsection{Flusso di Utilizzo da parte dello Studente}

\paragraph{Prenotazione.}
\begin{enumerate}
  \item Lo studente apre \file{/prenota} e sceglie la data, la durata e la sala;
        la pagina carica sale e posti con \file{GET /api/sale} e
        \file{GET /api/posti}.
  \item Sulla mappa sceglie un posto libero, eventualmente attiva il margine
        pendolare, e conferma.
  \item La pagina invia \file{POST /api/prenotazioni} con identificativo
        dell'utente, posto, data e orari.
  \item Il route handler controlla i campi obbligatori, legge utente, posto e
        sala, verifica che l'intervallo rientri nell'orario di apertura e che
        non esistano prenotazioni attive sovrapposte sul posto o per l'utente.
  \item Se i controlli sono superati crea la prenotazione nello stato
        \texttt{CONFERMATA}, registra l'evento \texttt{PRENOTAZIONE\_CREATA},
        invia una notifica e risponde 201; in caso di sovrapposizione risponde
        409.
\end{enumerate}

\diagramma[0.9\linewidth]{sequenza-prenotazione-as-is}{Sequence Diagram della prenotazione.}

\paragraph{Check-in.}
\begin{enumerate}
  \item Nel giorno della prenotazione la home mostra il QR code della
        prenotazione e il pulsante di conferma della presenza.
  \item La conferma invia \file{POST /api/prenotazioni/[id]/check-in}; il
        route handler verifica la sessione e che la prenotazione appartenga
        all'utente.
  \item Se il momento della richiesta cade nella finestra di quindici minuti
        attorno all'inizio, la prenotazione passa a \texttt{CHECK\_IN} e il
        posto a \texttt{OCCUPATO}; altrimenti la richiesta è rifiutata come
        scaduta.
\end{enumerate}

\diagramma[0.9\linewidth]{sequenza-checkin-as-is}{Sequence Diagram del check-in.}

\subsection{Flusso di Utilizzo da parte dello Staff}

\begin{enumerate}
  \item Il bibliotecario apre \file{/admin/scanner} e inquadra il QR code
        mostrato dallo studente.
  \item La pagina invia il contenuto del QR a
        \file{POST /api/admin/scanner/validate}, che verifica sessione e ruolo,
        la firma e la validità temporale del codice.
  \item Il route handler controlla che la prenotazione esista, appartenga allo
        studente indicato nel QR e non sia annullata, scaduta o già in corso;
        quindi registra il check-in, occupa il posto, crea l'evento e la
        notifica e mostra allo staff i dati dello studente.
  \item Dalla pagina \file{/admin/prenotazioni} lo staff può annullare una
        prenotazione (\file{POST /api/admin/prenotazioni}, azione
        \texttt{ANNULLA\_SINGOLA}): la prenotazione diventa
        \texttt{CANCELLATA}, il posto torna disponibile se era occupato e lo
        studente riceve una notifica.
\end{enumerate}

\diagramma[0.9\linewidth]{sequenza-staff-as-is}{Sequence Diagram dello staff.}

\subsection{Flusso delle Automazioni Pianificate}

\begin{enumerate}
  \item L'attività pianificata di Vercel (Vercel Cron) chiama
        \file{GET /api/cron/automations}, passando il segreto
        \texttt{CRON\_SECRET} nell'intestazione di autorizzazione.
  \item Il route handler confronta il token con il segreto configurato e
        invoca \texttt{runAllAutomations()}.
  \item Il servizio invia i promemoria di check-in alle prenotazioni che
        iniziano entro venti minuti, gli avvisi per i prestiti in scadenza e
        rilascia le prenotazioni \texttt{CONFERMATA} non confermate oltre il
        margine: stato \texttt{NO\_SHOW}, posto \texttt{DISPONIBILE}, notifica
        all'utente ed evento \texttt{NO\_SHOW\_AUTO}.
  \item Al termine restituisce un riepilogo delle operazioni eseguite.
\end{enumerate}

\diagramma[0.9\linewidth]{sequenza-automazioni-as-is}{Sequence Diagram delle automazioni pianificate.}

\end{document}
